\documentclass[11pt, a4paper, oneside]{article}
\usepackage[ngerman]{babel}
\usepackage{worksheet}

\begin{document}
	\author{L. Bung}
	\title{Wiederholung: Exponentialfunktionen}
	\subject{Mathematik}
	\class{FHR2}
	\maketitle
	
	\task{Funktionsgleichungen bestimmen}
	
	Bestimmen Sie die Funktionsgleichung der Exponentialfunktion $f(x) = a \cdot q^x$, die durch die angegebenen Punkte verläuft.
	
	\begin{enumerate}[label=\alph*)]
		\item $A(0 \mid 3)$ und $B(2 \mid 12)$
		\item $P_1(1 \mid 12)$ und $P_2(3 \mid 108)$
		\item $Q_1(-1 \mid 2)$ und $Q_2(2 \mid 54)$
	\end{enumerate}
	
	
	\task{Exponentialgleichungen lösen}
	
	Geben Sie die Lösung(en) der Exponentialgleichungen auf drei Nachkommastellen genau an.
	
	\begin{enumerate}[label=\alph*)]
		\item $5 \cdot 2^{x+1} = 80$
		\item $3^{2x-1} = 5^{x+1}$
		\item $4 \cdot e^{3x} - 12 \cdot e^{x} = 0$
		\item $7 \cdot 3^{x} + 2 = 65$
	\end{enumerate}
	
	
	\task{Basiswechsel}
	
	\begin{enumerate}[label=\alph*)]
		\item Schreiben Sie die Funktionsgleichungen in der Form $f(x) = a \cdot e^{k \cdot x}$.
		\begin{enumerate}[label=\roman*)]
			\item $f(x) = 5 \cdot 2^x$
			\item $g(x) = 12 \cdot 0{,}8^x$
		\end{enumerate}
		
		\item Schreiben Sie die Funktionsgleichungen in der Form $f(x) = a \cdot q^x$.
		\begin{enumerate}[label=\roman*)]
			\item $h(x) = 3 \cdot e^{1{,}2 \cdot x}$
			\item $k(x) = 20 \cdot e^{-0{,}5 \cdot x}$
		\end{enumerate}
	\end{enumerate}
	
	
	\task{Schnittpunkte von Exponentialfunktionen}
	
	Berechnen Sie den bzw. die Schnittpunkt(e) der beiden Exponentialfunktionen.
	\begin{enumerate}[label=\alph*)]
		\item $f(x) = 2 \cdot 3^{x+1}$ und $g(x) = 18 \cdot 2^{x-1}$
		\item $f(x) = e^x - 5$ und $g(x) = -6 \cdot e^{-x}$
	\end{enumerate}
	
\end{document}